\documentclass[12pt]{article}
\usepackage[T1]{fontenc}
\usepackage[utf8]{inputenc}
\usepackage{lmodern}
\usepackage{textcomp}
\usepackage{enumitem}
\usepackage{geometry}
\geometry{margin=1in}
\begin{document}
\begin{center}
Answer Key - Geography - Undergraduate Level Assessment 5 \
\small (Answers are ordered exactly as in the question paper.)
\end{center}

\section*{Multiple Choice}
\begin{enumerate}[label=\arabic*.]
\item \textbf{Answer:} B
\\ \textit{Options:} A) 12 years; B) 9 years; C) 6 years; D) 3 years
\\ \textit{Note:} The correct answer is 9 years because, as of 2020, global educational attainment data indicated that women aged 30 were generally less educated than their male counterparts, reflecting persistent gender disparities in access to education.
\item \textbf{Answer:} C
\\ \textit{Options:} A) 10\%; B) 20\%; C) 30\%; D) 40\%
\\ \textit{Note:} The correct answer is 30\% because, as of 2019, approximately 2.2 billion people worldwide lacked access to safely managed drinking water services, which translates to around 30\% of the global population.
\item \textbf{Answer:} A
\\ \textit{Options:} A) 9\%; B) 29\%; C) 59\%; D) 79\%
\\ \textit{Note:} The correct answer is 9\% because, according to various reports and data from international organizations, this figure accurately reflects the limited access to electricity in South Sudan due to ongoing conflict, lack of infrastructure, and economic challenges.
\item \textbf{Answer:} D
\\ \textit{Options:} A) shortly after World War II; B) in the 1960 s; C) around the time of World War I; D) in the early nineteenth century
\\ \textit{Note:} Most Latin American countries achieved independence in the early nineteenth century as a result of a series of revolutionary movements influenced by Enlightenment ideas and the decline of Spanish colonial power following the Napoleonic Wars.
\item \textbf{Answer:} B
\\ \textit{Options:} A) 36\%; B) 56\%; C) 76\%; D) 96\%
\\ \textit{Note:} The correct answer of 56\% reflects historical data indicating that, in 1950, a little over half of the global population could read and write, highlighting the significant impact of educational access and socio-economic factors on literacy rates during that period.
\end{enumerate}

\section*{True / False}
\begin{enumerate}[label=\arabic*.]
\setcounter{enumi}{5}
\item \textbf{Answer:} True
\\ \textit{Options:} A) True; B) False
\\ \textit{Note:} According to the passage: tripedal
\item \textbf{Answer:} False
\\ \textit{Options:} A) True; B) False
\\ \textit{Note:} The correct answer is: misguided agitation
\item \textbf{Answer:} False
\\ \textit{Options:} A) True; B) False
\\ \textit{Note:} The correct answer is: the Cortes or Junta
\item \textbf{Answer:} True
\\ \textit{Options:} A) True; B) False
\\ \textit{Note:} According to the passage: 1738
\item \textbf{Answer:} True
\\ \textit{Options:} A) True; B) False
\\ \textit{Note:} According to the passage: obligations established by agreement (express or implied) between private parties
\end{enumerate}

\section*{Long Form Response}
\begin{enumerate}[label=\arabic*.]
\setcounter{enumi}{10}
\item \textbf{Answer:} The Napoleonic Wars
\\ \textit{Note:} Expected answer: The Napoleonic Wars
\item \textbf{Answer:} 249,998
\\ \textit{Note:} Expected answer: 249,998
\end{enumerate}

\vfill
\noindent\textit{End of Answer Key}
\end{document}